% Fragmento listo para integrarse en el marco teórico.
% Las cuatro figuras se esperan en la ruta de figuras del documento principal:
% p1-poisson.pdf, p2-burgers.pdf, p3-allen-cahn.pdf y p4-dfg.pdf.
% Este fragmento interpreta los contratos de P1--P4; no sustituye la
% especificación experimental del capítulo de metodología.

\section{Lectura física de los problemas canónicos}
\label{sec:lectura-fisica-pdes}

Las ecuaciones diferenciales parciales del proyecto no son ejemplos
intercambiables. Cada una concentra una dificultad física distinta: Poisson
representa un equilibrio elíptico con variación distribuida; Burgers separa
transporte no lineal y difusión en una transición estrecha; Allen--Cahn
combina reacción, difusión y capas de interfaz; y el benchmark DFG exige
respetar simultáneamente inercia, viscosidad, incompresibilidad, geometría y
balances integrales. Esta lectura sirve para entender qué significa que una
aproximación represente bien un campo. Los contratos exactos del banco de
problemas se conservan en el capítulo de metodología; aquí se explican sus
consecuencias físicas y las magnitudes que hacen visibles.

Cada problema se acompaña de dos figuras con funciones distintas. La
ilustración en tinta, realizada con Fiziko, describe el sistema y sus datos.
La gráfica en color representa una solución exacta o numérica e identifica
su método de obtención. Los cálculos siguientes verifican referencias para
la explicación del problema; no evalúan todavía el mecanismo adaptativo
PIKAN propuesto.

\subsection{P1. Poisson: equilibrio y frecuencia distribuida}


\begin{figure}[htbp]\centering
\includegraphics[width=\linewidth,height=.30\textheight,keepaspectratio]{p1-problem.pdf}
\caption[P1: ilustración del problema de Poisson]{Ilustración de P1 con Fiziko: placa cuadrada, fuente distribuida y valor impuesto en los bordes. Puede interpretarse como una diferencia de temperatura respecto de la frontera. Para $k=4$, la fuente manufacturada alterna generación y extracción; el dibujo identifica los datos y no representa un campo calculado.}\label{fig:p1-problem}
\end{figure}


Considérese el cuadrado unitario $\Omega=(0,1)^2$ y el problema
\begin{equation}
-\Delta u=f_k \quad \text{en }\Omega,
\qquad u=0 \quad \text{en }\partial\Omega,
\qquad
u_k^*(x,y)=\sin(k\pi x)\sin(k\pi y),
\label{eq:expansion-p1-poisson}
\end{equation}
donde
\begin{equation}
f_k(x,y)=2(k\pi)^2u_k^*(x,y).
\label{eq:expansion-p1-source}
\end{equation}
La función $u_k^*$ se anula en los cuatro lados, por lo que satisface la
condición de Dirichlet sin corrección adicional. Además,
$\partial_{xx}u_k^*=-(k\pi)^2u_k^*$ y
$\partial_{yy}u_k^*=-(k\pi)^2u_k^*$; por tanto,
$-\Delta u_k^*=2(k\pi)^2u_k^*=f_k$. La igualdad es una verificación por
sustitución de la solución manufacturada, no una aproximación de los valores
de $u$.

La interpretación física más simple es un estado estacionario: el operador
$-\Delta$ mide cómo la curvatura del campo compensa una fuente distribuida.
Si $u$ representa temperatura en un medio homogéneo, $f$ puede representar
generación interna de calor; si representa un potencial, $f$ desempeña el
papel de una densidad de fuente. En ambos casos no hay una dirección de
propagación ni una condición inicial. La frontera fija el nivel del campo y
el interior ajusta su curvatura para satisfacer el balance.

El parámetro $k$ cambia la escala espacial sin cambiar la amplitud máxima de
$u_k^*$. Los ceros consecutivos de $\sin(k\pi x)$ están separados por
$1/k$, mientras que la curvatura y la fuente crecen como $k^2$. Para leer
la oscilación sin mezclarla con una región de interés arbitraria, puede
tomarse el corte
\begin{equation}
u_k^*\!\left(x,\frac{1}{2k}\right)=\sin(k\pi x),
\label{eq:expansion-p1-corte}
\end{equation}
porque $\sin(k\pi/(2k))=1$. Así, $k=1$ muestra una variación amplia y
$k=4$ muestra cuatro semiondas en la misma longitud. La familia permite
preguntar si la representación resuelve frecuencia distribuida en todo el
dominio; una ganancia en P1 no puede interpretarse automáticamente como
localización de capacidad en una interfaz o frente.

\paragraph{Solución y comprobación.}
Los mapas se obtienen de la fórmula exacta. Como comprobación independiente,
se discretiza $-\Delta$ con el stencil centrado de cinco puntos y se resuelve
el sistema disperso con Dirichlet homogéneo. Al pasar de 33 a 65 y 129 nodos
por eje, el error relativo discreto se reduce aproximadamente por cuatro en
cada refinamiento, coherente con segundo orden. El residual algebraico se
comprueba por separado: resolver bien el sistema discreto no elimina su
error de discretización. Esta comparación aplica el método de diferencias
finitas desarrollado más adelante sin usar datos de entrenamiento.
\begin{figure}[htbp]\centering
\includegraphics[width=\linewidth]{p1-solution.pdf}
\caption[P1: solución y comprobación]{Solución exacta de Poisson para $k=1$ y $k=4$, cortes en $y=1/(2k)$ y comprobación independiente por diferencias finitas para $k=4$. Se resuelve el sistema discreto con 33, 65 y 129 nodos por eje. El panel de error compara la aproximación con la fórmula exacta en los nodos interiores; la pendiente de referencia es de segundo orden.}\label{fig:expansion-p1-poisson}
\end{figure}


La forma débil hace explícita la misma lectura. Para una prueba $v$ que se
anula en $\partial\Omega$,
\begin{equation}
\int_\Omega \nabla u\cdot\nabla v\,\mathrm d x
 =\int_\Omega f_kv\,\mathrm d x.
\label{eq:expansion-p1-debil}
\end{equation}
El lado izquierdo mide el trabajo asociado a los gradientes y el derecho la
acción de la fuente. Una discretización puede aproximar esta identidad aun
cuando el campo local tenga error; por ello se deben distinguir balance
integral, error de campo y error de derivadas. Esta distinción es central al
comparar una aproximación neuronal con una referencia clásica
\cite{evans1998pde,fenics2016}.

\subsection{P2. Burgers: transporte no lineal, difusión y frente}


\begin{figure}[htbp]\centering
\includegraphics[width=\linewidth,height=.30\textheight,keepaspectratio]{p2-problem.pdf}
\caption[P2: ilustración del transporte y la difusión]{Ilustración de P2 con Fiziko: medio unidimensional en el que el propio campo controla el transporte y la viscosidad suaviza el frente. Las flechas explican el mecanismo de Burgers y las fronteras corresponden a trazas de la onda exacta. Es un modelo escalar, no una reconstrucción de un conducto tridimensional.}\label{fig:p2-problem}
\end{figure}


En $x\in[-1,1]$ y $t\in[0,1]$, el segundo contrato es
\begin{equation}
u_t+u\,u_x-\nu u_{xx}=0,
\qquad
u^*(x,t)=c-A\tanh q,
\qquad
q=\frac{A(x-ct-x_0)}{2\nu}.
\label{eq:expansion-p2-burgers}
\end{equation}
La condición inicial y las dos condiciones de frontera son las trazas de
$u^*$:
\begin{equation}
u(x,0)=u^*(x,0),
\qquad
u(-1,t)=u^*(-1,t),
\qquad
u(1,t)=u^*(1,t).
\label{eq:expansion-p2-conditions}
\end{equation}
Se fijan $A=0.5$, $x_0=-0.25$,
$c\in\{0,0.5\}$ y $\nu\in\{0.02,0.005\}$. El centro del frente está en
$x_f(t)=ct+x_0$: con $c=0$ permanece estacionario y con $c=0.5$ recorre
media unidad durante el intervalo temporal. La frontera dependiente del
tiempo no es un forzamiento añadido; es la traza necesaria para mantener la
onda exacta al entrar o salir del dominio.

La sustitución muestra qué términos se equilibran. Si
$S=\operatorname{sech}^2q$, entonces
\begin{equation}
u_x=-\frac{A^2}{2\nu}S,
\qquad
u_t=-c\,u_x,
\qquad
u_{xx}=\frac{A^3}{2\nu^2}S\tanh q.
\label{eq:expansion-p2-derivatives}
\end{equation}
Como $u-c=-A\tanh q$,
\begin{equation}
u_t+u u_x=(u-c)u_x
 =\frac{A^3}{2\nu}S\tanh q
 =\nu u_{xx}.
\label{eq:expansion-p2-balance}
\end{equation}
El término $u u_x$ transporta la forma mediante la velocidad local; la
difusión $\nu u_{xx}$ impide que la transición se vuelva discontinua. El
ancho característico $\ell=2\nu/A$ vale $0.08$ cuando $\nu=0.02$ y
$0.02$ cuando $\nu=0.005$. El máximo del gradiente tiene magnitud
$A^2/(2\nu)$, que pasa de $6.25$ a $25$ al reducir $\nu$. La amplitud del
campo apenas cambia, pero la derivada que entra en el residual se concentra
en una región cuatro veces más estrecha.

Esta relación explica por qué P2 es una prueba de resolución localizada. Una
aproximación puede tener valores razonables fuera del frente y, al mismo
tiempo, perder la curvatura que equilibra transporte y difusión. También
explica por qué movimiento y estrechamiento son dificultades distintas: $c$
cambia la posición de la señal en el espacio-tiempo, mientras $\nu$ cambia
la escala de sus derivadas. La familia y sus ondas viajeras tienen antecedente
analítico en la literatura de Burgers \cite{burgerswaves2025}; los cuatro
cruces de parámetros pertenecen al contrato finito de este proyecto.

\paragraph{Solución y comprobación.}
La gráfica representa la onda exacta y sus trazas, no un ajuste a muestras.
La sustitución de sus derivadas verifica la identidad de Burgers para los
cuatro cruces de $c$ y $\nu$ del contrato, con residual del orden del
redondeo de doble precisión. Por ello esta fórmula sirve como referencia
directa para medir error de campo y error del frente en futuras aproximaciones.
\begin{figure}[htbp]\centering
\includegraphics[width=\linewidth]{p2-solution.pdf}
\caption[P2: solución y comprobación]{Solución exacta de Burgers con $A=0.5$, $c=0.5$, $x_0=-0.25$ y $\nu=0.02$: mapa espacio--tiempo y perfiles a $t=0,0.5,1$. La línea discontinua del mapa sigue $x_f(t)=ct+x_0$. Las condiciones iniciales y de frontera son las trazas de la misma solución.}\label{fig:expansion-p2-burgers}
\end{figure}


\subsection{P3. Allen--Cahn: reacción, energía e interfaz}


\begin{figure}[htbp]\centering
\includegraphics[width=\linewidth,height=.30\textheight,keepaspectratio]{p3-problem.pdf}
\caption[P3: ilustración de fases e interfaz]{Ilustración de P3 con Fiziko: regiones de fase y una zona de transición, con difusión y reacción. Los símbolos $-1$ y $+1$ identifican estados preferidos. La disposición de fases es conceptual; la evolución del dato inicial se muestra por separado en la gráfica de solución.}\label{fig:p3-problem}
\end{figure}


El tercer problema conserva el contrato del ejemplo documentado de DeepXDE:
\begin{align}
u_t-d u_{xx}-5(u-u^3)&=0,
&&x\in[-1,1],\quad t\in[0,1],\quad d=0.001,\\
u(x,0)&=x^2\cos(\pi x),
&&u(-1,t)=u(1,t)=-1.
\label{eq:expansion-p3-allen-cahn}
\end{align}
Puede escribirse como un flujo de gradiente de la energía
\begin{equation}
\mathcal E[u]=\int_{-1}^{1}
\left[\frac d2|u_x|^2+W(u)\right]\,\mathrm d x,
\qquad
W(u)=\frac54(u^2-1)^2.
\label{eq:expansion-p3-energy}
\end{equation}
En efecto, $W'(u)=5(u^3-u)$ y la ecuación equivale a
$u_t=d u_{xx}-W'(u)$. La reacción $5(u-u^3)$ empuja el campo hacia los
estados estables $u=\pm1$ y hace inestable el estado $u=0$; el término
difusivo penaliza cambios espaciales demasiado bruscos. La solución debe
resolver ambos impulsos, además de acomodarse al dato inicial y a las
fronteras fijadas en $-1$.

La escala de una interfaz se obtiene al equilibrar difusión y reacción. Para
el perfil estacionario idealizado que conecta $-1$ con $+1$,
\begin{equation}
u_{\rm kink}(x)=
\tanh\!\left(\frac{x-x_c}{\ell_{\rm AC}}\right),
\qquad
\ell_{\rm AC}=\sqrt{\frac{2d}{5}}=0.02.
\label{eq:expansion-p3-kink}
\end{equation}
La cifra $0.02$ es la escala del perfil de transición para $d=0.001$; no es
una afirmación de que la evolución con el dato inicial del contrato tenga una
única interfaz de ese ancho en todo tiempo. El estado inicial
$x^2\cos(\pi x)$ puede producir varias zonas de transición, y las
condiciones $u(\pm1,t)=-1$ influyen en su evolución. La máscara física de
una interfaz debe construirse a partir del campo de referencia cuando se
evalúa el problema, no a partir de la forma idealizada de esta explicación.

\paragraph{Solución numérica y sensibilidad a la resolución.}
Se discretiza el término difusivo con diferencias centradas de segundo
orden y se integra el sistema de ecuaciones ordinarias mediante BDF, usando
su Jacobiano disperso analítico. Las fronteras se fijan exactamente en $-1$.
Se comparan 401, 801 y 1601 nodos espaciales sobre 301 tiempos de salida
comunes. La diferencia máxima entre los dos últimos refinamientos es
$4.36\times10^{-4}$, y el cociente entre diferencias sucesivas es 3.94.
Al reducir las tolerancias relativas de $10^{-9}$ a $10^{-11}$ y absolutas
de $10^{-11}$ a $10^{-13}$ en la malla fina, la diferencia máxima es
$6.14\times10^{-9}$: en esta comparación domina la sensibilidad espacial.
Estas son comprobaciones de autoconvergencia y de tolerancia, no una cota
demostrada del error absoluto.

La energía discreta usa diferencias entre nodos para aproximar
$\int d|u_x|^2/2+W(u)\,dx$ y decrece en los tiempos de salida. Para una
solución suave con fronteras fijas, multiplicar
$u_t=d u_{xx}-W'(u)$ por $u_t$ e integrar da
\[
\frac{\mathrm d\mathcal E}{\mathrm dt}=-\int_{-1}^1|u_t|^2\,dx\leq0.
\]
Así se relaciona el panel de energía con el mecanismo físico. Allen--Cahn
describe un parámetro de orden no conservado: la reacción puede cambiar su
integral espacial, aunque la energía disminuya. La forma de los frentes
resulta del problema inicial completo, a diferencia del perfil ideal de una
sola interfaz en la recta.
\begin{figure}[htbp]\centering
\includegraphics[width=\linewidth]{p3-solution.pdf}
\caption[P3: solución y comprobación]{Solución numérica del problema Allen--Cahn con $d=0.001$, $u(x,0)=x^2\cos(\pi x)$ y $u(\pm1,t)=-1$. Se emplean diferencias centradas y BDF, con 1601 nodos en la solución mostrada. Los paneles incluyen campo, cortes, energía discreta y diferencias entre refinamientos, evaluadas en nodos y tiempos comunes.}\label{fig:expansion-p3-allen-cahn}
\end{figure}


Allen--Cahn introduce una diferencia importante respecto de Burgers. En
Burgers la transición se mueve por un balance entre advección y difusión; en
Allen--Cahn la transición también puede cambiar porque el campo desciende por
un potencial local con dos estados preferidos. Por eso el valor de $u$ y su
gradiente tienen interpretaciones simultáneas: estar cerca de $\pm1$ indica
fase, mientras que $|u_x|$ identifica una interfaz. Una representación que
suaviza demasiado la interfaz puede reducir una medida global y aun así
desplazar el lugar donde el campo cambia de fase. La formulación y las
condiciones se mantienen coherentes con el ejemplo de DeepXDE versión 1.12.2
\cite{deepxdeAC}.

\subsection{P4. Navier--Stokes: flujo DFG alrededor de un cilindro}


\begin{figure}[htbp]\centering
\includegraphics[width=\linewidth,height=.30\textheight,keepaspectratio]{p4-problem.pdf}
\caption[P4: ilustración del canal con cilindro]{Ilustración de P4 con Fiziko: canal y cilindro fijo, entrada parabólica, paredes sin deslizamiento y salida natural. Las flechas de entrada son datos prescritos. La geometría y las fronteras se distinguen de las líneas de corriente calculadas en la figura de solución.}\label{fig:p4-problem}
\end{figure}


El benchmark DFG 2D-1 considera flujo estacionario incompresible alrededor de
un cilindro. El dominio es el canal
\begin{equation}
\Omega=([0,2.2]\times[0,0.41])
\setminus\overline{B_{0.05}(0.2,0.2)},
\label{eq:expansion-p4-domain}
\end{equation}
con densidad unitaria, viscosidad cinemática $\nu=0.001$ y campos
$\mathbf u=(u,v)$ y $p$. El sistema local es
\begin{align}
u u_x+v u_y+p_x-\nu(u_{xx}+u_{yy})&=0,\\
u v_x+v v_y+p_y-\nu(v_{xx}+v_{yy})&=0,\\
u_x+v_y&=0.
\label{eq:expansion-p4-ns}
\end{align}
Los dos primeros renglones expresan el balance de momento: la convección
transporta momento, el gradiente de presión lo redistribuye y la viscosidad
disipa gradientes. El tercero impone que el volumen no se comprima; no es una
ecuación auxiliar que pueda omitirse sin cambiar el problema.

La entrada del canal prescribe el perfil parabólico
\begin{equation}
\mathbf u_{\rm in}(y)=
\left(\frac{4(0.3)y(0.41-y)}{0.41^2},0\right),
\label{eq:expansion-p4-inlet}
\end{equation}
cuya velocidad máxima es $0.3$ en la línea media y cuya media transversal
es $0.2$. Las paredes del canal y la circunferencia del cilindro tienen
condición de no deslizamiento, $\mathbf u=\mathbf0$. En la salida se usa,
con la convención del benchmark para el operador Laplaciano,
\begin{equation}
\nu\,\partial_n\mathbf u-p\mathbf n=\mathbf0.
\label{eq:expansion-p4-outlet}
\end{equation}
Esta condición deja salir el flujo sin fijar artificialmente un valor de
presión adicional. Con $L=0.1$ como diámetro y $U=0.2$ como velocidad de
referencia, $Re=UL/\nu=20$, de modo que el caso es laminar y permite
examinar la estela sin convertir el problema en un estudio de turbulencia.

La forma integral conecta la PDE con las cantidades de ingeniería. Como
$\nabla\cdot\mathbf u=0$ y $\nu$ es constante, para un volumen de control
$K\subset\Omega$ sin fuerza volumétrica se obtiene
\begin{equation}
\int_{\partial K}
\left[\mathbf u(\mathbf u\cdot\mathbf n)+p\mathbf n
      -\nu(\nabla\mathbf u)\mathbf n\right]\,\mathrm d s=\mathbf0.
\label{eq:expansion-p4-integral}
\end{equation}
Cada cara intercambia flujo convectivo, presión y esfuerzo viscoso. Sobre el
cilindro, la tracción según la convención usada en el contrato es
$\boldsymbol\sigma_b\mathbf n$, con
$\boldsymbol\sigma_b=\nu\nabla\mathbf u-pI$. Por tanto,
\begin{equation}
\begin{pmatrix}F_D\\F_L\end{pmatrix}
 =\int_S\boldsymbol\sigma_b\mathbf n\,\mathrm d s,
\qquad
C_D=\frac{2F_D}{U^2L},
\qquad
C_L=\frac{2F_L}{U^2L}.
\label{eq:expansion-p4-forces}
\end{equation}
En esta integral, $\mathbf n$ apunta hacia fuera del cilindro y hacia el
fluido; es el opuesto de la normal exterior del dominio fluido sobre esa
frontera interna. Declarar esta orientación evita invertir los signos de
las fuerzas.

La solución permite leer el balance espacialmente. El obstáculo desvía el
flujo hacia los pasos laterales, donde aumenta la rapidez, mientras que el
no deslizamiento obliga a anularla en la superficie. La presión participa
en esa desviación y en el frenado frontal; la viscosidad transmite el efecto
de las paredes al fluido cercano. La presión no se calcula mediante una
ley de estado de compresibilidad: en este modelo ajusta el campo de
velocidad para satisfacer simultáneamente el balance y la restricción
$\nabla\!\cdot\mathbf u=0$. Las líneas de corriente son tangentes a la
velocidad calculada; al ser estacionario este caso, también describen
trayectorias de partículas ideales del campo resuelto.

El benchmark publica valores de arrastre, sustentación y caída de presión que
se usan como controles externos en la metodología; no son resultados
producidos por esta propuesta \cite{dfgRe20}.

\paragraph{Solución FEM y comprobación.}
La discretización empleó elementos triangulares Taylor--Hood \(P_2/P_1\) para
velocidad y presión. Cada nivel se generó con Delaunay y una frontera
poligonal explícita para el cilindro; el número de segmentos de esa frontera
creció junto con la resolución de fondo. La convección se continuó mediante
Picard y se remató con correcciones de Newton. La tabla resume la corrida
completa que produjo los campos de la figura.

\begin{center}
\small
\begin{tabular}{lrrrrrr}
\hline
nivel & $n_x\times n_y$ & segmentos & elementos & $N_u$ & $N_p$ & Picard+Newton\\
\hline
grueso & $34\times8$ & 32 & 734 & 3168 & 425 & 18+1\\
medio  & $52\times12$ & 48 & 1526 & 6456 & 851 & 18+1\\
fino   & $76\times18$ & 64 & 3100 & 12904 & 1676 & 18+1\\
\hline
\end{tabular}
\end{center}

En el nivel fino, el residuo débil libre fue
\(2.11\times10^{-18}\), el residuo de continuidad débil
\(1.95\times10^{-18}\), el error de traza de entrada
\(4.14\times10^{-17}\) y el máximo de no deslizamiento
\(7.22\times10^{-16}\). El cierre integral de flujo fue
\(2.54\times10^{-15}\) relativo, con flujos de entrada y salida
\(-0.0820\) y \(0.0820\), respectivamente; la traza de salida evaluada con
la condición natural tuvo norma \(7.18\times10^{-7}\). La norma \(L^2\) de
la divergencia evaluada dentro de los elementos bajó de
\(1.40\times10^{-1}\) a \(6.21\times10^{-2}\); la restricción débil de
continuidad se reporta por separado porque es la que impone el espacio mixto.

Los funcionales obtenidos en las tres mallas fueron
\[
\begin{array}{c|rrr}
 & C_D & C_L & \Delta p\\\hline
\text{gruesa} & 5.51699 & 0.00544 & 0.11364\\
\text{media}  & 5.52807 & 0.01015 & 0.12186\\
\text{fina}   & 5.52153 & -0.00651 & 0.11555
\end{array}
\]
con \(\Delta p=p(0.15,0.2)-p(0.25,0.2)\) y
\(\boldsymbol\sigma=\nu\nabla\mathbf u-pI\) integrada con la normal
exterior al cilindro. Como control externo, el benchmark publica
\(C_D=5.57953523384\), \(C_L=0.010618948146\) y
\(\Delta p=0.11752016697\)~\cite{dfgRe20}. La comparación entre niveles da
 cambios relativos medio--fino de \(0.24\%\) en velocidad y \(0.99\%\) en
 presión de media cero, pero \(5.18\%\) en \(\Delta p\); además, el signo de
 \(C_L\) no se estabiliza. Por ello estos resultados documentan una solución
 FEM reproducible y sus balances, pero no se presentan como certificación de
 convergencia de todas las cantidades integrales. Las imágenes de velocidad,
 presión y líneas de corriente se obtuvieron exclusivamente interpolando el
 campo resuelto y enmascarando el mismo polígono sólido de la malla.

\begin{figure}[htbp]\centering
\includegraphics[width=\linewidth]{p4-solution.pdf}
\caption[P4: solución FEM de velocidad y presión]{Solución FEM calculada del problema estacionario de Navier--Stokes
DFG 2D-1 a $\mathrm{Re}=20$ en la malla fina Taylor--Hood $P_2/P_1$.
Las tres filas muestran, respectivamente, la magnitud de velocidad, la
presión y las líneas de corriente calculadas a partir de los grados de
libertad resueltos; el cilindro poligonal y su máscara sólida se conservan en
los tres paneles. La entrada tiene velocidad máxima $0.3\,\mathrm{m/s}$ y
media $0.2\,\mathrm{m/s}$, las paredes y el cilindro son no deslizantes, y la
salida usa $\nu\,\partial_n\mathbf u-p\mathbf n=0$. La corrida satisface
residuo débil y balance de flujo, pero la tabla de refinamiento no estabiliza
el funcional de sustentación; por tanto la imagen se presenta como solución
reproducible con verificación de balances, sin certificado de convergencia
integral.}\label{fig:expansion-p4-dfg}
\end{figure}


El carácter canónico de DFG proviene de que un campo visualmente plausible no
garantiza el balance. Una pequeña violación de $u_x+v_y=0$ acumula flujo
espurio; una presión con nivel o signo inconsistente altera la tracción; y
una salida mal impuesta puede contaminar la estela. Por eso el desempeño debe
leerse en campos y en observables integrales: arrastre, sustentación y caída
de presión. La geometría curva y el no deslizamiento junto al cilindro
introducen gradientes donde una descripción global suave puede parecer
adecuada lejos de la superficie y fallar en la región que determina las
fuerzas.

\subsection{Lectura conjunta del cuarteto}

P1 cambia la frecuencia en todo el dominio; P2 concentra una derivada en un
frente que puede moverse; P3 combina un potencial local con una interfaz cuya
escala se obtiene de un balance reacción--difusión; y P4 exige compatibilidad
entre campos, frontera curva y cantidades integrales. La secuencia permite
separar cuatro preguntas físicas: si una representación resuelve variación
distribuida, si conserva una transición transportada, si mantiene estados y
interfaces, y si satisface un balance vectorial con incompresibilidad.

Las parejas de ilustración y solución acompañan esas preguntas desde dos
perspectivas: los datos del sistema y el campo que los satisface. En P1 y P2
se dispone de soluciones exactas; en P3 y P4 se usan aproximaciones numéricas
con comprobaciones de resolución y balance. Su propósito es preparar la lectura de
las métricas y referencias del capítulo de metodología. La evidencia de una
hipótesis sobre adaptación solo podrá provenir de comparaciones reproducibles
con esas referencias; las ilustraciones presentes aquí explican el problema,
pero no sustituyen tal evidencia.
